\section{Introduction}
\label{sec:intro}

Many real-world classification problems carry an externally specified
hierarchy over class labels: artistic-style movements branch into one
another \citep{tan2019wikiart}, biological taxa nest into clades, and
ontologies of products, diseases, and documents are organised as
trees. A natural way to fold that hierarchy into a probabilistic
classification model is to add a tree-structured regularizer to the
layout of class-conditional parameters. We study the simplest version
of this: a prototype classifier whose class-conditional likelihood is
a softmax over distance-to-prototype, with a regularizer on the
matrix of pairwise prototype distances pulling it toward the
tree-distance matrix. The fitted classifier is then a balance between
a class-likelihood term (paintings of style $k$ should be close to
prototype $\mathbf{p}_k$) and a tree-shape term (prototypes
themselves should respect the hierarchy).

The choice of latent manifold for the prototypes is the question this
paper investigates. Forcing a tree's exponentially-growing leaf count
into Euclidean $\mathbb{R}^d$, whose ball volume grows only
polynomially in radius, induces unavoidable distortion of pairwise
distances; this distortion has been characterised analytically and
shown to drop sharply on the Poincar\'e ball $\mathbb{B}^d_c$, whose
own exponential volume growth matches a tree's
\citep{nickel2017poincare,sala2018representation}. In our model, the
tree-structured regularizer is a function of pairwise prototype
distances and should therefore be cheaper to satisfy on the
hyperbolic manifold than on the Euclidean one without distorting the
class-conditional likelihood. Whether that geometric advantage
realises in practice on a real hierarchical-classification problem is
what we test.

We instantiate the question on WikiArt-Refined \citep{tan2019wikiart},
a 27-style 81{,}446-painting corpus where the class hierarchy
(Renaissance branches into Baroque, Baroque into Rococo, Rococo into
Romanticism, on into Impressionism and Cubism) is part of the domain
and was hand-built from standard art-history references. Frozen
CLIP ViT-B/16 features \citep{radford2021clip} feed an MLP head and a
prototype classifier; the manifold of the classifier output is the
only design choice that varies between Euclidean and hyperbolic runs.
The same tree-structured regularizer (a normalised squared-error
pull on the prototype distance matrix) applies in both geometries.

We sweep regularized maximum-likelihood fits across 150
seed-replicated configurations spanning embedding dimension
$d \in \{2, 4, 8, 16, 32, 64\}$, curvature $c \in \{0.1, 0.3, 1, 3\}$,
regularizer strength $\lambda \in \{0, 0.1, 0.3, 1, 3\}$, and three
reference-tree variants (default lineage, chronological era grouping,
flat null). Three empirical findings emerge.

\paragraph{Local hierarchy.} Hyperbolic prototypes recover local tree
structure substantially better than matched Euclidean ones. Sibling
recall@5 at $d{=}8$ is $0.195$ versus $0.149$ on the default tree;
aggregated over 18 paired seed configurations the mean gap is
$+8.7$\,pp on sibling recall and $+15.2$\,pp on cousin recall, with
sign agreement $0.94$ and paired-$t$ $p<10^{-4}$. This effect is
preserved when sibling sets are redefined from data-driven empirical
trees built either from CLIP or DINOv2 \citep{oquab2024dinov2}
features, and across the regularizer-strength sweep.

\paragraph{Calibration.} A logistic-regression baseline on raw CLIP
features achieves $64.1\%$ top-1, statistically tied with the
Euclidean fit's $64.3\%$. A $k$-NN baseline on the same features
achieves sibling recall@5 of $0.160$, statistically tied with the
Euclidean fit's $0.149$ but $3.5$\,pp below the hyperbolic fit's
$0.195$. Read against these baselines, the hyperbolic fit is the only
one of our trained models that improves the encoder on local
hierarchy; the Euclidean fit matches the encoder on classification
but adds no detectable structural value over it.

\paragraph{Global tree fidelity.} Mean tree distortion against the
default tree favours Euclidean significantly, but
prototype-tree Spearman is not significantly different from zero
across the full sweep ($p=0.68$), and the small effect that does
appear flips sign between empirical reference trees built from CLIP
and from DINOv2 feature centroids. We therefore do not claim a
geometry winner on global tree fidelity.

The clean claim our experiments support is narrower than ``geometry
matters globally'': hyperbolic latent structure adds local-retrieval
value over the encoder while Euclidean latent structure does not, and
that local advantage is robust to which of three reference trees
defines ``local.''

\paragraph{Related work.} Hyperbolic prototype embeddings are an
established tool for tree-shaped data
\citep{nickel2017poincare,sala2018representation,ganea2018hyperbolic}.
\citet{khrulkov2020hyperbolic} report a similar asymmetry on
few-shot benchmarks (helps neighborhood structure, leaves top-1
unchanged); we confirm that pattern on a hand-built taxonomy with a
tree-aware regularizer, and quantify it against linear and $k$-NN
calibration baselines that previous reports of the same asymmetry do
not include. Tree-distance-matrix penalties over learnable points are
used implicitly in hyperbolic hierarchical clustering
\citep{chami2020trees}; here we use one as an explicit, tunable
regularizer.
